%% LyX 2.4.4 created this file.  For more info, see https://www.lyx.org/.
%% Do not edit unless you really know what you are doing.
\documentclass[english]{article}
\usepackage[T1]{fontenc}
\usepackage[latin9]{inputenc}
\usepackage{enumitem}
\usepackage{amsmath}
\usepackage{amsthm}
\usepackage{amssymb}
\usepackage{geometry}
\geometry{verbose,tmargin=1in,bmargin=1in,lmargin=1in,rmargin=1in}

\makeatletter
%%%%%%%%%%%%%%%%%%%%%%%%%%%%%% Textclass specific LaTeX commands.
\numberwithin{equation}{section}
\numberwithin{figure}{section}
\newlength{\lyxlabelwidth}      % auxiliary length 
\theoremstyle{definition}
\newtheorem*{defn*}{\protect\definitionname}

%%%%%%%%%%%%%%%%%%%%%%%%%%%%%% User specified LaTeX commands.
\usepackage{fancyhdr}
\usepackage{lastpage}
\pagestyle{fancy}
\usepackage{enumitem}
\usepackage{ifthen}
\everymath=\expandafter{\the\everymath\displaystyle}
%\DeclareMathSizes{10}{10}{10}{10}
\usepackage{multicol}

\usepackage{tikz}
\usetikzlibrary{patterns}
\usetikzlibrary{plotmarks}
\usepackage{pgfplots}
\pgfplotsset{compat=newest}

\usepackage{calc}
\usepackage[nomessages]{fp}

\usepackage{datetime}
% Added by lyx2lyx
\setlength{\parskip}{\smallskipamount}
\setlength{\parindent}{0pt}

\makeatother

\usepackage{babel}
\providecommand{\definitionname}{Definition}

\begin{document}
\renewcommand{\author}{Dr.\,James Ashe}

\newcommand{\classNum}{335}

\newcommand{\classSec}{A}

\newcommand{\nameType}{Name: \hspace{-4cm}}

\newcommand{\examType}{Homework 5}

\newcommand{\Date}{Due date: \formatdate{26}{4}{2013}} %%\AdvanceDate[12] \today

%\newdate{my_date}{\day}{\month}{\year}%   
%\AdvanceDate[-5] 
%\displaydate{my_date}

%%First page's header%%%%%%%%%%%%%%%%%%%%%%
\fancypagestyle{firststyle} %%header settings for first pagr
{
	\fancyhead[L]{MTH \classNum{}(\classSec{}) \author{}\\
	\textbf{\examType{}} \Date{}}
	\fancyhead[C]{\textbf{\nameType}}
	\setlength{\headsep}{14pt}
}
%%%%%%%%%%%%%%%%%%%%%%%%%%%%%%%%%%%%%%%%%%

%%Next page's header%%%%%%%%%%%%%%%%%%%%%%%
%\setlist{nolistsep} %eliminates space between list items
\lhead{MTH \classNum{}(\classSec{})} 
\chead{\textbf{\nameType}} 
\cfoot{\thepage{}~of~\pageref{LastPage}} 
\setlength{\headsep}{5pt} 
%%%%%%%%%%%%%%%%%%%%%%%%%%%%%%%%%%%%%%%%%%%

%%Various settings; see the preamble for other settings%%%%%
%%\setenumerate[0]{leftmargin=12pt,itemindent=0pt,labelwidth=0pt}
\renewcommand{\labelenumi}{\textbf{\arabic{enumi}.}}
\renewcommand{\labelenumii}{\textbf{\alph{enumii}.}}
\thispagestyle{firststyle}
%%%%%%%%%%%%%%%%%%%%%%%%%%%%%%%%%%%%%%%%%%

\textbf{Homework Problems. }\emph{Solutions should be written/typed
in numerical order and clearly labeled. Unsolved problems should be
included but left blank.}
\begin{enumerate}
\item Let $\left(\mathbb{Q},+\right)$ be the group of rational numbers
under addition. 

\begin{enumerate}
\item Is $\mathbb{Z}$ a subgroup of $\left(\mathbb{Q},+\right)$?
\item Is $\mathbb{Z}^{+}=\left\{ 0,1,2,3,\dots\right\} $ a subgroup of
$\left(\mathbb{Q},+\right)$?\vfill
\end{enumerate}
\item Determine if the following are subgroups.

\begin{enumerate}
\item Is $\mathbb{Q}$ a subgroup of the group, $\left(\mathbb{C},+\right)$,
the complex numbers under addition?
\item Is $\mathbb{Z}-\left\{ 0\right\} $ a subgroup of the group, $\left(\mathbb{Q}-\left\{ 0\right\} ,\cdot\right)$,
the non-zero rational numbers under multiplication?\vfill
\end{enumerate}
\item Prove that the intersection of two subgroups $H$ and $K$ of a group
$G$ is a subgroup.\vfill
\item Let $H_{1},H_{2},\dots$ be subgroups of $G$. Prove that $H_{1}\cap H_{2}\cap\cdots$
is a subgroup of $G$.\vfill
\item Recall that $\mathbb{Z}_{6}=\left\{ 0,1,2,3,4,5\right\} $ is a group
under modular addition. 

\begin{enumerate}
\item Show that $H=\left\{ 0,2,4\right\} $ and $K=\left\{ 0,3\right\} $
are subgroups of $\mathbb{Z}_{6}$.
\item Show that $H\cup K$ is \emph{not }a subgroup of $\mathbb{Z}_{6}$.
\item Is the union of two subgroups of $G$ necessarily a subgroup $G$?\vfill
\end{enumerate}
\item Let $G$ be an abelian group and $H$ be a subgroup of $G$. Prove
that $S(H)=\left\{ x|x\in G\mbox{ and }x\cdot x\in H\right\} $ is
also a subgroup of $G$.\vfill
\item Let $G$ be a group. Prove that $G$ has a subgroup with one element.
This group is called the \emph{trivial group}.\vfill
\end{enumerate}
\begin{defn*}
Let $G$ be a group. If there is a single $g\in G$ such that every
element in $G$ is equal to $g^{n}$ for some $n\in\mbox{\ensuremath{\mathbb{N}}}$
then $G$ is a \emph{cyclic }group and \emph{x }is called the \emph{generator}
of $G$.
\end{defn*}
\begin{enumerate}[resume]
\item Prove that:

\begin{enumerate}
\item $\mathbb{Z}_{n}$ is a cyclic group under modular addition for any
$n\in\mathbb{Z}$. (you can use the fact that it is a group).
\item $\mathbb{Z}$ is a cyclic group. \vfill
\end{enumerate}
\item Find all generators of 

\begin{enumerate}
\item $\mathbb{Z}_{6}$, $\mathbb{Z}_{8}$, and $\mathbb{Z}_{10}$ (note:
$\mathbb{Z}_{n}$ is often notated $Z_{n}$).\vfill
\end{enumerate}
\end{enumerate}
\begin{defn*}
Let $X=\left\{ 1,2,\dots,n\right\} $. The \emph{symmetric group on
X, $S_{n}$, }is the collection of all bijections $\sigma:X\rightarrow X$.
That is, it is the set of all functions that rearrange the elements
of $X$. ($S_{n}$ is also called the \emph{permutation group of the
set X}).
\end{defn*}
\begin{enumerate}[resume]
\item Let $S_{3}$ be the symmetric group of permutations on three elements.
Show that $S_{3}$ is:

\begin{enumerate}
\item Non-abelian.\vfill
\item Not cyclic.\vfill
\end{enumerate}
\end{enumerate}

\end{document}
